\documentclass[12pt]{article}
\usepackage[utf8]{inputenc}
\usepackage[T1]{fontenc}
\usepackage{hyperref}
\title{Botolan Sambal Language Resources for K-12 Curriculum Development}
\author{Compiled by Claude Sonnet 5.5 (claude-sonnet-5-5) for Andr\'as Kornai}
\date{}
\begin{document}
\maketitle

\section{Introduction}
Botolan Sambal (ISO 639-3 sbl, Glottocode boto1242) is an Austronesian language of the Sambalic group, spoken mainly in the municipalities of Botolan and Cabangan in Zambales, Philippines. SIL (2000) gives about 32,900 speakers. Speakers of the related Sambal (Tina) reportedly understand about 70 percent of Botolan, so the two should not be treated as the same language. Search found a New Testament and a dramatized audio recording, but no dictionary or Botolan-specific school materials.

\section{Primary Educational Resources}
No Botolan-specific primer or textbook turned up. DepEd's mother tongue policy added Sambal to the languages used for Grade 1 instruction from school year 2013--2014. The search results do not say whether Sambal there means the Botolan variety, so I cannot confirm that. A Zambales schools division page lists a training workshop on radio-based scripts for MTB-MLE ``Sambal G1'', which shows some teacher-training activity for Sambal in the province.
\begin{itemize}
\item DepEd, mother tongue order adding Sambal for Grade 1. \url{https://deped.gov.ph/?p=7545}
\item DepEd Zambales, radio-based script workshop for MTB-MLE (Sambal G1). \url{https://zambales.deped.gov.ph/sept-1-2-schools-division-online-training-workshop-on-the-development-of-radio-based-scripts-in-english-and-mother-tongue-based-multi-lingual-education-mtb-mle-in-g2-and-sambal-g1}
\end{itemize}

\section{Linguistic and Reference Materials}
No Botolan dictionary or grammar turned up. Reference pages on the Sambalic languages and the UP Diliman linguistics language capsule on Sambal are general background. The UP page also says OLAC lists Sambal texts and language descriptions, which could be searched further, though those may cover other varieties.
\begin{itemize}
\item Wikipedia, Botolan language (ISO 639:sbl). \url{https://wikipedia.com/wiki/ISO_639:sbl}
\item Wikipedia, Sambalic languages. \url{https://en.wikipedia.org/wiki/Sambalic_languages}
\item UP Diliman Linguistics, Sambal language capsule. \url{https://linguistics.upd.edu.ph/the-katig-collective/language-capsules/sambal}
\item Omniglot, Botolan. \url{https://www.omniglot.com/writing/botolan.htm}
\end{itemize}

\section{Religious and Community Resources}
A Botolan Sambal New Testament with Psalms and Proverbs, ``Hay halit\'a nin Diyos hay kawkanta boy kawkahalita-an'', was published in 2005 by Wycliffe Bible Translators (Bible League print edition recorded in the Princeton catalogue). The eBible.org page marks it as all rights reserved, so reuse needs permission.
\begin{itemize}
\item eBible.org, Botolan Sambal New Testament. \url{https://ebible.org/bible/details.php?id=sbl}
\item eBible.org, copyright and details. \url{https://ebible.org/sbl/copyright.htm}
\item Scripture Earth, Botolan Sambal. \url{https://scriptureearth.org/00eng.php?iso=sbl}
\item Princeton University Library catalogue record. \url{https://catalog.princeton.edu/catalog/SCSB-6191955}
\end{itemize}

\section{Digital and Audio Resources}
MegaVoice hosts a dramatized audio New Testament in Botolan Sambal. Scripture Earth also points to Bible.com, Bible.is and a Gospel of Mark film.
\begin{itemize}
\item MegaVoice, New Testament (Dramatized), Sambal, Botolan. \url{https://megavoice.com/media-cloud/m032602-new-testament-dramatized-sbl-sambal-botolan-audio-bible/}
\end{itemize}

\section{Curriculum Development Context}
Official MTB-MLE policy covers Sambal in Grade 1, but I found no Botolan-specific learner materials, and it is unclear how far Botolan is actually used in classrooms. The best native-language text sources are the New Testament and its audio, which are copyrighted. A curriculum effort would need to commission original materials and confirm with the Zambales division office which variety its Sambal materials use. This dossier is short because little was found.
\end{document}